\documentclass[conference]{IEEEtran}
\usepackage{amsmath,amssymb,booktabs,array,url,xcolor}
\usepackage[hidelinks]{hyperref}
\newcommand{\stopgrad}{\operatorname{sg}}
\title{Adaptive Video Preprocessing Techniques for Optimizing Video Coding for Machines (VCM) on NVIDIA Jetson Orin NX}
\author{\IEEEauthorblockN{Anonymous Author(s)}\\\IEEEauthorblockA{Identities pending human approval}}

\begin{document}
\maketitle
\begin{center}
\fcolorbox{red}{red!7}{\parbox{0.92\columnwidth}{\centering\bfseries PRE-RESULTS MANUSCRIPT DRAFT --- DO NOT SUBMIT}}
\end{center}

\begin{abstract}
Video coding for machines must preserve features required by downstream analysis while respecting bitrate, human-viewing quality, and edge-compute constraints. We present a standards-compatible framework that places trainable QP-conditioned neural pre- and postprocessors around a frozen H.264/AVC or H.265/HEVC codec. Real decoded pixels and measured bitstream BPP are used in every forward pass, while a frozen predictive-entropy proxy supplies only the backward Jacobian. The objective combines rate, a frozen task loss, DINOv2 feature consistency, and human-oriented spatial, temporal, perceptual, and adaptive-DCT terms. The wrappers export independently to TensorRT around the hardware codec on NVIDIA Jetson Orin NX. This paper defines the method and locked evaluation protocol; no new empirical result is claimed because the required dataset, checkpoint, and target hardware were unavailable at manuscript generation time.
\end{abstract}

\begin{IEEEkeywords}
video coding for machines, neural preprocessing, standard codec, straight-through estimation, DINOv2, FiLM, Jetson Orin NX
\end{IEEEkeywords}

\section{Introduction}
Conventional H.264/AVC and H.265/HEVC codecs are standardized and widely hardware accelerated, but their native objectives do not directly optimize a downstream neural task \cite{wiegand2003avc,sullivan2012hevc}. Replacing them with learned codecs may sacrifice compatibility. Standard-compatible preprocessing retains deployed bitstreams while adapting their inputs.

Lu et al. learn a quantization-adaptive image preprocessor and identify the benefit of real-codec forward values with proxy backpropagation \cite{lu2024preprocessing}. Sandwiched Compression jointly trains neural wrappers around a standard codec \cite{guleryuz2024sandwiched}. Zhao et al. introduce video-specific task-aware preprocessing and a virtual codec \cite{zhao2025video}; RPP contributes rate-perception and adaptive-DCT ideas \cite{ma2023rpp}. We combine these directions with DINOv2 semantic features \cite{oquab2023dinov2}, FiLM QP conditioning \cite{perez2018film}, and local Video Swin attention \cite{liu2022video}.

Our contributions are: (i) identity-initialized QP-FiLM Video Swin Lite and FiLM-3D wrappers; (ii) real H.264/H.265 forward values and proxy-only gradients for both pixels and rate; (iii) a supervised, DINOv2, perceptual, temporal, and adaptive-DCT objective; and (iv) a deployment and measurement contract for Jetson Orin NX.

\section{Method}
For source clip $x$ and quantization parameter $q$, the preprocessor produces $z=P_\theta(x,q)$. The frozen codec gives $y=C_q(z)$ and elementary-stream rate $r$. The postprocessor outputs $\hat{x}=R_\phi(y,q)$. The task analyzer consumes $\hat{x}$ by default; an ablation consumes $y$.

The standard codec is non-differentiable. With frozen reconstruction and rate proxies $(\hat C_\psi,\hat r_\psi)$, training uses
\begin{align}
y_{\mathrm{ST}} &= \hat C_\psi(z,q)+\stopgrad(C_q(z)-\hat C_\psi(z,q)),\\
r_{\mathrm{ST}} &= \hat r_\psi(z,q)+\stopgrad(r_q(z)-\hat r_\psi(z,q)).
\end{align}
The forward values are therefore exact real-codec outputs while gradients with respect to $z$ come from the proxy. This estimator remains approximate and is accepted only after reconstruction/rate audits and a real-codec gradient-direction probe.

$P_\theta$ uses a spatial Conv3D patch embedding and alternating regular/shifted 3-D windows without temporal downsampling. $R_\phi$ is a shallow 3-D residual U-Net with depthwise-separable blocks and one spatial down/up stage. QP embeddings modulate intermediate features through FiLM. Zero-initialized RGB residual heads make both wrappers exact identities initially.

The loss is
\begin{equation}
\mathcal L=w_r\frac{r_{\mathrm{ST}}}{r_{\mathrm{anchor}}(q)}+w_t\mathcal L_{\mathrm{task}}+w_s\mathcal L_{\mathrm{DINO}}+w_h\mathcal L_{\mathrm{human}}.
\end{equation}
$\mathcal L_{\mathrm{DINO}}$ is cosine distance between frozen source/restored CLS and patch tokens. The human term combines Charbonnier, a compact three-scale MS-SSIM surrogate, optional LPIPS \cite{zhang2018lpips}, adjacent-frame gradient consistency, and adaptive block-DCT sparsification. The DCT term suppresses weak high-frequency coefficients while not directly penalizing strong edge/motion coefficients.

\section{Experimental Protocol}
The implemented primary task is Kinetics-400 action recognition \cite{kay2017kinetics}. The locked protocol uses 16 frames at stride 2 and $128\times128$ spatial size. Scientific evaluation uses real FFmpeg \texttt{libx264} and \texttt{libx265} at QP $\{30,32,35,37,40,42,45\}$. Every video/QP produces paired anchor, pre-only, and sandwich rows. BPP is elementary-stream bits divided by $THW$.

We report Top-1/Top-5, PSNR, externally validated MS-SSIM, LPIPS, and VMAF where available, plus task and perceptual BD-rate. Undefined BD-rates are not extrapolated. Required ablations remove postprocessing, DINO, adaptive DCT, or QP-FiLM; replace Video Swin with CNN; move the task tap before postprocessing; and replace real-forward training with proxy-forward training. Three seeds and paired video-level bootstrap intervals are required.

On Jetson Orin NX, pre/post models are separate ONNX/TensorRT engines around NVIDIA hardware encoder/decoder elements. The report records module variant, carrier, JetPack/L4T, CUDA, cuDNN, TensorRT, precision, power mode, clocks, codec configuration, and hashes. At least 30 warm-up and 200 timed iterations produce component/end-to-end p50/p95 latency, throughput, peak memory, VDD\_IN power, and energy/frame.

\section{Results Status}
No new empirical result is available in the supplied workspace. Table~\ref{tab:results} is intentionally a result schema. Prior-work values---including the $-20.3\%$ and $-14.6\%$ image-codec ablation of Lu et al.---must not populate this table.

\begin{table}[t]
\centering
\caption{Primary result schema. NM means not measured.}
\label{tab:results}
\begin{tabular}{llccc}
\toprule
Codec & Method & Task BD-R & PSNR BD-R & MS-SSIM BD-R\\
\midrule
H.264 & Anchor & Ref. & Ref. & Ref.\\
H.264 & Pre-only & NM & NM & NM\\
H.264 & Sandwich & NM & NM & NM\\
H.265 & Anchor & Ref. & Ref. & Ref.\\
H.265 & Pre-only & NM & NM & NM\\
H.265 & Sandwich & NM & NM & NM\\
\bottomrule
\end{tabular}
\end{table}

\section{Discussion and Limitations}
Exact forward values remove a major train--deployment mismatch but do not make the proxy Jacobian correct, motivating a gradient-direction audit. Joint postprocessing can improve human reconstruction while shifting the analyzer input; reporting task taps before and after restoration exposes this trade-off. DINOv2 may reduce analyzer overfitting but is framewise, and adaptive DCT may discard subtle task cues. Both require ablation.

This release has not been trained or measured on Kinetics-400 or Jetson Orin NX. Action recognition is implemented; tracking is not. TensorRT compatibility depends on the target release. Task-aware compression can remove information important to people, future tasks, safety review, or forensics, so deployment requires retention, privacy, and per-class failure policies.

\section{Conclusion}
We define an auditable standard-codec video sandwich for joint machine, human, and edge-device objectives. The implementation is ready for the empirical stage, while its claims remain deliberately pre-results until the locked experiments and human review are complete.

\bibliographystyle{IEEEtran}
\bibliography{references}
\end{document}
